\subsection{Challenges} \label{sec:challenges}
Throughout this thesis, we have seen that the field of AI for code has made significant progress
over the years. However, state-of-the-art models still struggle with software engineering tasks, especially
at larger scopes and higher levels of logical complexity. Next, we discuss
several key challenges in the field, motivated by the insights we have learned.
The challenges below are not meant to be an exhaustive list, but rather a synthesis of
several key themes inspired from the work in this thesis.

\subsubsection{Evaluation and Benchmarks} \label{sec:c-evaluation}

In Chapter~\ref{chap:right-metrics}, we saw the importance of focusing on the right metric and eliminating biases in evaluation. This lesson is the inspiration of the
first challenge, which discusses three familiar difficulties in benchmarking AI coding systems.

\begin{aicodegreenbox}
Today's code LLM evaluations focus on a narrow set of tasks, suffer from potential contamination, and do not reliably measure real-world software engineering abilities.
\\
\newline
\textit{Potential solutions}: \ref{sec:d-data}
\end{aicodegreenbox}

Our taxonomy of tasks and measures highlights some of the shortcomings of today's evaluations and benchmarks. For example, the majority of today's coding evaluations have no level of human intervention, with a few, such as Copilot-Arena \citep{chi2025copilotarena}, having low to medium autonomy. 
% \ms{nit: copilot-arena style evals}
HumanEval, MBPP, APPS, CodeContests, and LiveCodeBench are all at function-level scope, with low to medium-high logical complexity. 
Commit0 \citep{zhao2024commit0}, SWE-Bench \citep{jimenez2024swebench}, TestGenEval~\citep{jain2024testgeneval}, RefactorBench~\citep{gautam2024refactorbench}, SWE-Lancer \citep{miserendino2025swe} are at project-level scope with low to medium logical complexity. %Today, we lack benchmarks 1) at project-level scope with high logical complexity, 2) for many tasks other than code generation, and 3) requiring human intervention.

\textbf{Task Diversity and Capability Isolation}: 
Current coding evaluations primarily focus on the code generation task, while most of the tasks discussed in Section~\ref{sec:tasks-milestones} are either not studied such as Code QA or only studied in limited scopes like EvalPerf~\citep{evalperf}, vulnerability detection~\citep{mei2024arvo}, formal verification~\citep{sun2024clover}. 
As more agent-based approaches are introduced for software engineering (e.g. pairing a code generation model with a debugging model), these engineering-related capabilities beyond just code generation will be important towards designing a maximally performant system. 
Solely relying on end-to-end coding evaluations that focus on the overall correctness of a codebase makes it difficult to precisely measure progress and learn from the failure modes on individual tasks.


\textbf{Contamination}: 
Data contamination is a serious issue that, if not taken into account, can affect the soundness of various conclusions drawn from a set of benchmark results. In coding, the performance of LLMs on competitive programming \citep{xu2024benchmark, jain2024livecodebenchholisticcontaminationfree} and SWE-Bench \citep{aleithan2024swe} tasks has been shown to degrade over time, indicating the possibility of older problems being contaminated due to public exposure on the internet. For simpler function-level HumanEval style problems, \citet{matton2024leakage} suggest three potential causes of contamination: direct data leakage (benchmarks are on GitHub), synthetic data leakage (there are only a limited number of interview problems), and overfitting to test sets (benchmark hacking). In addition, for code, contamination can be hard to detect, as semantically equivalent code that is syntactically distinct could be thought of as contamination \citep{riddell2024quantifying}. A recent benchmark, the Konwinski Prize (\url{https://www.kprize.ai/}), is a promising way to fairly evaluate SoTA LLM models by only using new GitHub issues.

\textbf{Construct Validity}:
Construct validity refers to how closely a measurement reflects the underlying concept. 
Given the implications of rapid performance improvement in the AI for the code domain, it is essential to have high-construct validity benchmarks evaluating how well programming agents can perform.
While benchmarks like SWE-Bench come close, user experiences do not currently match rapid performance gains obtained from them. 
This is partially because many desiderata in software engineering cannot be described cleanly via automated unit testing. Things like multi-turn code generation, designing an UI, and writing clean and idiomatic code are all difficult to quantitatively measure with precision. Designing reliable proxy metrics for these desired goals remains a challenge. 

% \todo{most of SWE cannot be described via automated unit testing WD}

% Existing works that focus on 

% \todo{talk about challenge in programing systems community to measure the success of new PLs. notorously dificult to eavluate. User studies fail to capture maintainability, learning effects.}



%\todo{I wonder if it makes sense to make a list of benchmarks and which category they belong to}

% Given the rapid pace of progress of LLMs, 

% Machine learning progress is in a large part guided by making steady progress (``hill climbing'') on some set of benchmarks.
% For example, \naman{ImageNet served ...}.
% Similarly, AI for code has witnessed a rising interest in the development of better coding evaluations

% Thus the choice of benchmarks developed and used in this process 

% Machine learning approaches are typically about optimization and progress is made through improving (``hill climbing'') on benchmarks. 
% Therefore, developing appropriate benchmarks that allow generalization to 
% the choice of benchmarks developed or used during the machine learning life-cycle is very important.
% Here, we draw upon learnings from prior works and present directions for 

% The rate of progress in artificial intelligence is considerably aff


\subsubsection{Effective Tool Usage} \label{sec:c-tool-use}

In our work on LeanDojo and ProofOptimizer, we saw that tools like the Lean compiler formed the basis for proof search and RL. Similarly, benchmarks like SlopCodeBench and VibeCodeBench revealed that models needed to use execution feedback and tests in order to succeed at longer-horizon tasks. These examples point towards the same broader lesson: systems must be able to effectively work with tools.

\begin{aicodegreenbox}
While software engineers use a wide suite of programming tools when programming, most of today's AI coding systems do not invoke tools. AI needs to be able to select which tool to use, decide how to use it, and interpret the outputs in order to continue making progress on the task.
\newline \\
\textit{Potential solutions}: \ref{d:sec-agents-tools}
\end{aicodegreenbox}

Software engineering has witnessed the development of various open and proprietary tooling support for programming, debugging, analysis, and code management over time.
% Software engineering is a long-standing field and has support for mature tooling support f
% which has led to development of various mature tooling to support developers in programming and debugging tasks. 
For example, program analysis tools provide static and dynamic assurances on code correctness. 
Print statements and debuggers are used for dynamically analyzing and debugging programs at a fine-grained level.
Beyond programming, such tools are richly integrated into the entire software development lifecycle, e.g., code navigation or search, reviewing code, CI testing.

There have been efforts combining LLMs with tools such as calculators and search engines \citep{schick2023toolformer, patil2023gorilla}. 
However, effective integration of LLMs with software engineering tools is a more challenging problem. 
Several early works have incorporated such tool feedback in code generation in an automated fashion, for example, linter or execution feedback in \citep{olausson2023self, zhong2024ldb, gehring2024rlef}.
However, these works do not actively \textit{interact} with tools. 
More recently, programming agents have started incorporating tool use within their workflows termed as Agent-Computer-Interface~\citep{yang2024sweagent}. 
These tools range from aiding in general search (\texttt{grep}), providing code editor for making changes \citep{openhands,claude35swebench}, language server for static analysis \citep{liu2024marscode},
dependency analyzer \citep{bairi2024codeplan}, terminal access for bash commands including code execution ~\citep{yang2024sweagent}, debugger \citep{bigsleep}. 
% For example, ~\citet{sonnet35swe, openhands} uses string replace code editor, some of these tools at increasingly rapid pace, but their use is relatively limited \citep{openhands, yang2024sweagent}.

% \todo{Cite \citep{deshpande2024class}}

% \todo{More SWE-Agent stuff? [Naman]}

\textbf{Dynamic and Effective Tool Usage}: While many efforts combine LLMs and agents with tools, they do not achieve fully dynamic and effective software engineering tool usage. This involves an AI system seamlessly and proactively integrating appropriate tools depending on the task at hand. There are a few challenges towards achieving this goal. First, the AI system must identify which tools could potentially be useful for the task at hand. Second, the system then needs to decide when to invoke the tool. A complex debugging task might require the use of \texttt{pdb} or \texttt{gdb} to track intermediate program states, while looking at input-output pairs may be sufficient for simple debugging tasks. Third, the agent then must figure out how to invoke the tool. If the agent knows that a certain function in a program has an error, it may wish to walk through only that function instead of the entire code from start to finish. Finally, the agent needs to incorporate the output provided by the tool in order to inform its next steps, e.g. edit the code if a bug was uncovered or run the tool again otherwise.

% \textit{Example: Performance Instrumentation}: CSI \citep{schardl2017csi} is a way to instrument software by programatically inserting hooks to track objects such as memory loads/stores and function entry/exits to measure code coverage and performance. For an LLM agent to use CSI effectively, it must 1) learn the CSI API, 2) decide when and how to use CSI (such as when it writes a suspected performance bottleneck), and 3) use the CSI output to decide how to proceed.

% \textbf{Leveraging program analysis}: Apart from LLMs, there are a wide variety of program analysis tools designed to prevent bugs and ensure code correctness properties. Abstract interpretation \citep{cousot1977abstract} is a technique to compute over-approximations of program state in order to prove the soundness of program properties at points in the code. Concolic testing \citep{godefroid2005dart, sen2005cute} finds bugs in software by combining concrete and symbolic execution. Model checking \citep{clarke1997model} is a way to prove properties of execution traces via temporal logic. Linting and type-checking \citep{cardelli1996type} provide a static check to ensure that variables, expressions, and functions adhere to a programming language's rules. 


% There are a lot of programming and debugging tools that people use that most AI programming systems probably don't leverage or know how to use effectively:

% \begin{itemize}
%     \item Print statement debugging.
%     \item Debuggers (pdb/gdb).
%     \item Linters.
%     \item For optimization, things like Valgrind.
% \end{itemize}

\subsubsection{Long-Horizon Code Planning} \label{sec:c-long-horizon}

One of the key lessons from SlopCodeBench was that as models worked on longer-horizon tasks, their code quality often degraded. Similarly, in ProofOptimizer, we observed that models could produce formally correct code that was unreadable or difficult to maintain. The next challenge summarizes this broader issue: AI coding agents need to perform long-horizon planning, not merely produce locally correct artifacts.

\begin{aicodegreenbox}
Large software engineering projects often require long-term planning about the design and structure of the code. LLMs struggle at two key aspects of this: designing good, lasting abstractions and respecting modularity and code quality principles.
\newline \\
\textit{Potential solutions}: \ref{sec:d-rl}
\end{aicodegreenbox}

When working on large projects, engineers and tech leads often make complex decisions about how to design and structure the code to best support the various functionalities that will eventually be needed. To build a long-lasting software system, an engineer must know the potential paths that the system's evolution might take. This requires domain expertise and experience in how different code structures require different forms of extension. We believe that today's language models are unable to perform this level of sophisticated planning.

\textbf{Designing Good Abstractions}:
One instance of long-horizon code planning is choosing the right abstractions from the outset. An API designed with good abstractions will allow new features to be implemented seamlessly with minimal user overhead, while an API designed with poor abstractions may lead to excessive code duplication, refactoring, or even debugging. We discuss two examples of this, library learning and data representation.

\textit{Library learning}: Designing APIs and libraries are designed with useful abstractions often leads to more code reuse and more intuitive interfaces. The challenge of library learning is to derive a library of useful abstractions from a corpus of programs by abstracting out common reusable features \citep{ellis2021dreamcoder,10.5555/3692070.3693967,10.1145/3571234}. While the traditional library learning literature has focused primarily on code reuse, a truly effective library must also prioritize ease of use and maintainability, as well as be robust and adaptable to future extensions. %within the software system.

\textit{Data representation}: The choice between data structures leads to a variety of trade-offs when it comes to performance aspects such as memory usage and processing speed. For example, database engineers need to decide between various data models, storage formats, and indexing methods to balance performance. 

\textbf{Modularity and Code Quality}: %\todo{written code may not be of high quality, lack modularity}
LLMs are trained and optimized primarily for code correctness with insufficient focus on other aspects of code like quality and maintainability. This is further exacerbated with large scale reinforcement learning being performed using test cases which can lead to unintended consequences regarding code quality, as correct but poor quality code is still often given a high reward.
Empirically, it has been observed that LLM written solutions are often more complex than human-written counterparts. 
% \todo{cite? personally observed this a bit in competition coding where even o1 solutions are wayyy more complex. WD: can we cite the paper "Library Learning Doesn’t: The Curious Case of the Single-Use Library" saying library reuse is not trivial for LLM? \url{https://arxiv.org/pdf/2410.20274}}
For example, \citet{jain2024r2e} identified that LLMs prefer to repeat existing code instead of making use of abstractions in the existing code. 
One aspect of code quality is modularity, ensuring that code does not get duplicated too often. Here, \citet{berlot2024library} identified that library or tool reuse is non-trivial for LLMs in coding and formal math domains. 

% \todo{Comments from Jiawei
% \\
% - “Modularity and Code Quality” …. code quality is pretty important (to me at least) and could have a lot of connection with many sections in the manuscript:
% \\
% - Reward: R1-Zero like approaches tend to produce ugly code as code quality is not modelled as reward — it’s also hard to be modelled — except for using static analyzers to screen CWE issues such as CodeGuru, CodeQL, Bandit — these signals diversifies the reward (we have wip projects on this).
% --> {naman - empirical}
% \\
% - Also code quality might not be planned but post-processed — so I feel this is a bit more related to code review/editing/transformation.
% }
% \citep{kang2024revisiting}


% \textbf{Code Debugging}:
% \naman{
% Bug-fixing can be particularly challenging since the bugs may present themselves in considerably higher-order effects of the underlying issue.
% For example, \todo{find a swebench or github issue amplifying this aspect}.
% Solving bugs from such bug reports requires reasoning and codebase understanding to localize the cause of such underlying failure. \todo{this should be well-studied in SE research, cite from there?}
% Next, fixing the issue requires rethinking the implementation to fix the issue. 
% These implementation changes should be carefully planned, attempting to come up with \textit{most general} version of the changes while also keeping in mind other parts of the codebase that are affected.
% }


% \textbf{Abstractions in Code Refactoring}: Being able to perfectly predict the evolution of a software project is nearly impossible, and even highly experienced senior engineers with domain expertise fail to design all the abstractions correctly on the first try. Therefore, software often has to undergo major global refactorings. Fundamentally, refactoring is a process of mapping components in the old abstraction to those in the new one in a way that maintains semantic equivalence. This can be difficult for language models because these mappings can be relatively complex: large sections of code can be consolidated or split, the relationships between components may be redefined, and successfully maintaining the semantics properly requires precision. These challenges are reminiscent of those faced by the programming languages community when designing good refactoring tools.
% \naman{hmm i think refactoring is "too concrete" to talk about planning/abstractions -- the abstract bit in refactoring is making design decisions keeping long-term in mind..
% the content here on maintaining mapping seems lower level changes
% }
% First, there is a specification issue, as writing a full specification for a refactoring can be extremely complex, if not impossible. Second, there is a long-context issue, as refactorings can involve 

% In a sense, refactoring is about swapping out one set of abstractions for a new set of abstractions that is easier to use or can better support new use cases. Refactoring can be difficult for LLMs for several reasons. 

%  Finally, refactoring requires a deep understanding of both the original set of abstractions and the new ones. 
 
% \todo{refactoring. highly experienced senior engineers with domain expertise fail to get the abstractions right on the first try. a lot of things have to go through major global refactorings. discuss some of the challenges that are hard to achieve through pure LLM. 1) lots of LOC, 10k+, context window; 2) describing new structure and how it relates to the old structure. this is one of the challenges of PL tools too. 3) hard to express refactoring with a simple specification. MSR agentic refactoring work.  }

% \todo{Other points about long-horizon code planning?}

% \naman{\textbf{Tradeoffs}: I think talking about tradeoffs should be useful here.. rich space with long-term feedback..
% also acknowledging humans don't get it right in first attempt and major software do refactor
% /VLLM just did a major refactoring for v1 release actually!!}

\subsubsection{Semantic Understanding of Codebases} \label{sec:c-global-understanding}

Both CRUXEval and the Counterfeit Conundrum revealed that models were weak at understanding code, especially subtly incorrect code. Although today's models can now understand these function-level snippets, it remains important that they understand large codebases at scale, especially as
more and more software is being produced by AI.

\begin{aicodegreenbox}
Being able to effectively write code relies on having a strong semantic understanding (somewhat like a world model) of the entire codebase: structurally seeing how various parts of the code go together, knowing what is implemented where, understanding how the algorithms work, and keeping track of program invariants at certain program points. LLMs struggle with this global semantic understanding
\newline \\
\textit{Potential solutions}: \ref{sec:d-data}
\end{aicodegreenbox}

% \naman{why is this not a part of abstractions? -- comment based on subsec title}

% After working with a codebase for a prolonged time, programmers have a global understanding and mental model of the codebase. This includes overall codebase structure, different algorithms, stylistic aspects, data representation, program invariants, package versions, test coverage, and the interplay between different functions. This holistic understanding is necessary for performing many tasks.

A global and holistic semantic understanding of a codebase is important for performing almost all code-related tasks. For example, let's say an engineer is asked to improve the runtime performance of a query. To do so, they must first understand the codebase's structure well enough to know where all the pieces of the algorithm are implemented. Then, they need to understand the algorithm and implementation in detail. This includes both the high-level algorithm (including its time complexity) and the low-level implementation details to identify both algorithmic and implementation bottlenecks. Finally, after coming up with a solution, they must then return to their understanding of the global code structure so that they can integrate their new algorithm without introducing new bugs. 

% \textbf{Understanding invariants for code refactoring}: In order to refactor a large codebase, it is essential to have a complete understanding of how different parts of the codebase interact. These interactions are governed by a set of rules known as program invariants, properties about the code that are assumed to hold. While some of these invariants can be maintained and guaranteed explicitly with assertions, many of these invariants are implicit and may only be obvious to developers working on the codebase itself. When performing code refactoring, developers often have a mental model of these invariants and know what needs to be changed so that these invariants continue to hold. LLMs, on the other hand, lack such a nuanced understanding of programs and may not be able to infer or maintain complex invariants.

% \todo{Example of LLM not being able to understand an invariant.}

% \textbf{Understanding program execution for code debugging}: Debugging code is another task that relies on having a complete global understanding of a codebase. In order to debug programs productively, we need a good internal model of the intermediate execution states of a program. Step-by-step debuggers are generally helpful because the programmer has a mental model of properties that should be true at certain program states, and the debugger can surface whether those properties are indeed true. 

LLMs struggle at semantic understanding of codebases for several reasons. First, the way that code is pieced together can be relatively intricate, and understanding all these complex relationships can be difficult. Second, code can often have units with high logical complexity that contain custom algorithms that may never have appeared anywhere in the training data. Finally, because a disproportionately large number of LLM training tokens are spent on code generation rather than other coding tasks, models may lack a holistic awareness and world model of code. 

One desiderata is that models can generalize knowledge across various coding tasks \citep{roychoudhury2025will}. However, this may not be straightforward as just training on more tasks: \citet{gu2024cruxeval} found that coding models fine-tuned on additional natural language/code pairs saw significant improvements on code generation but did not transfer to improving code understanding and execution reasoning. While there have been successful efforts to augment code LLM training with execution information to improve general coding capabilities \citep{ni2024next, ding2024traced}, imbuing code LLMs with a general and holistic understanding of code remains an important challenge today.

% \todo{Example of LLM not being to reason about the execution of a piece of code}

% \todo{\url{https://www.expresscomputer.in/news/quantifying-github-copilots-impact-on-code-quality-ai/104480/} cite and also r2e discussion section}


\subsubsection{High Logical Complexity and OOD Domains} \label{sec:c-ood-domains}

In Chapter~\ref{chap:measure-shapes}, we saw how a benchmark can shape the direction of the research that follows.
Throughout the thesis, we saw a number of benchmarks that were more and more difficult as models improved in their capabilities.
This challenge is about the next frontier of these challenges, for which most of today's benchmarks do not address: high logical complexity and OOD domains.

\begin{aicodegreenbox}
Tasks such as writing highly concurrent code or discovering  performance optimizations have a high logical complexity, often proving difficult for even the best human coders. Similar to solving research-level math problems, these out-of-distribution domains are very hard for LLMs.
\newline \\
\textit{Potential solutions}: \ref{sec:d-rl}
\end{aicodegreenbox}

Some programming tasks are challenging for even the best human programmers, requiring approaches with a very high logical complexity. Examples of tasks that fall into this category include superoptimizing programs, discovering attacks for purportedly secure code, writing performant compilers, optimizing GPU kernels \citep{ouyang2024kernelbench}, and writing very error-prone and very technical code. 

% \todo{Flesh out the difficulties more}

% \todo{Add more examples of superhuman / OOD domains}

\textbf{Limits of Symbolic Techniques}: When it comes to applying symbolic techniques to these tasks, there are a few limiting factors that make them difficult to tackle. First, for synthesis-style tasks, the search space can be very large. Deductive and rewrite-based synthesis techniques are unable to explore a majority of the search space. Second, verifiers can be limited in power, such as when dealing with properties in concurrency or weak memory models. Third, many domains lack clean models to reason about properties, such as dealing with memory bandwidth in GPU kernels. 

Because they are hard for humans, these tasks are very rarely in the training data of today's language models. They have unique, domain-specific, challenges that making generalizing from existing data difficult. For these problems, language models rely heavily on feedback-driven search algorithms \citep{mankowitz2023faster}, and it can be difficult to navigate the search space effectively. In addition, many of these tasks lack feedback mechanisms, which is crucial for AI to pick up learning signals. When designing a complex algorithm or data structure, it is often hard to know if you are on the right track until you get to the correct result. When writing code for a large multithreaded operation, it may be hard to know if the algorithm has concurrency issues until all the parts are fully fleshed out. Without feedback, incremental improvement is nearly impossible.
